%% Solide soumis à DEUX forces : même support (l'axe de la barre), même intensité,
%% sens opposés. Vers l'extérieur = traction ; vers l'intérieur = compression.
\documentclass[tikz,border=4pt]{standalone}
\usepackage{liaisons}
\begin{document}
\begin{tikzpicture}[x=1cm,y=1cm,
  force/.style={-{Stealth[length=6pt]}, line width=1.2pt, solideA},
  support/.style={line width=0.5pt, dash pattern=on 3pt off 2pt, black!55}]

% --- macro locale : la barre inclinée à 45°, percée à ses deux extrémités ---
\newcommand\barre{%
  \begin{scope}[rotate=45]
    \draw[support] (-2.3,0) -- (2.3,0);
    \draw[line width=1.1pt, solideB, fill=solideB!10, rounded corners=6pt]
          (-1.1,-0.17) rectangle (1.1,0.17);
    \draw[line width=0.9pt, solideB, fill=white] (-0.85,0) circle (0.1);
    \draw[line width=0.9pt, solideB, fill=white] ( 0.85,0) circle (0.1);
  \end{scope}}

%% ------------------------- (a) traction -----------------------------
\begin{scope}[shift={(0,0)}]
  \barre
  \begin{scope}[rotate=45]
    \draw[force] (1.2,0) -- (2.25,0) node[anchor=south west, inner sep=1pt, font=\small, text=solideA] {$\vec{F}$};
    \draw[force] (-1.2,0) -- (-2.25,0) node[anchor=north east, inner sep=1pt, font=\small, text=solideA] {$\vec{F}$};
  \end{scope}
  \node[font=\small, anchor=north, align=center] at (0,-2.05) {(a) \textbf{traction}\\ flèches vers l'extérieur};
\end{scope}

%% --------------------------- « ou » ---------------------------------
\node[font=\small] at (2.5,0) {ou};

%% ----------------------- (b) compression ----------------------------
\begin{scope}[shift={(5.0,0)}]
  \barre
  \begin{scope}[rotate=45]
    \draw[force] (2.25,0) -- (1.2,0) node[anchor=north west, inner sep=1pt, font=\small, text=solideA] {$\vec{F}$};
    \draw[force] (-2.25,0) -- (-1.2,0) node[anchor=south east, inner sep=1pt, font=\small, text=solideA] {$\vec{F}$};
  \end{scope}
  \node[font=\small, anchor=north, align=center] at (0,-2.05) {(b) \textbf{compression}\\ flèches vers l'intérieur};
\end{scope}

%% ---------------------------- bilan ---------------------------------
\node[draw=black!45, fill=black!3, rounded corners=3pt, inner sep=5pt, anchor=north,
      font=\small] at (2.5,-2.9) {même support, même intensité, sens opposés};
\end{tikzpicture}
\end{document}
